\documentclass[11pt]{article}

\usepackage[T1]{fontenc}
\usepackage{lmodern}
\usepackage{amsmath,amssymb,amsthm}
\usepackage[margin=1in]{geometry}
\usepackage{microtype}
\usepackage[hidelinks,pdftitle={A Four-Job Counterexample to an Almost-Exact Fathoming Rule for Job Sequencing and Tool Switching},pdfauthor={Xiaojian Zeng}]{hyperref}

\newtheorem{theorem}{Theorem}
\newtheorem{remark}{Remark}
\newcommand{\cT}{\mathcal T}
\newcommand{\cM}{\mathcal M}

\title{A Four-Job Counterexample to an Almost-Exact Fathoming Rule\\
for Job Sequencing and Tool Switching}
\author{Xiaojian Zeng\\
\small PBC School of Finance, Tsinghua University, Beijing 100190, China\\
\small \texttt{xiaojianzeng@foxmail.com}}
\date{September 2026}

\begin{document}
\maketitle

\begin{abstract}
In 2026, Locatelli, C\^ot\'e, and Coelho introduced an almost-exact
fathoming rule for an insertion-based search for the job sequencing and tool
switching problem.  Their five-job example shows that this rule can fathom a
branch containing a globally optimal completion.  We disprove the proposed
lower boundary of five jobs.  An instance with four jobs, four tools, and
magazine capacity three has, under either orientation of its uniquely selected
root pair, a reachable proper node that is fathomed although it has a globally
optimal completion.  The root pair and next job are unique, both relevant
KTNS magazine histories are forced, and the competing completion-pruning rule
does not apply.  The proof is an explicit certificate; exact enumeration is
used only as a reproducibility check.
\end{abstract}

\section{Setting and contribution}

Let the jobs be $J=\{1,\ldots,n\}$ and the tools be
$\cT=\{1,\ldots,m\}$.  Job $j$ requires a nonempty set
$T_j\subseteq\cT$, and a magazine configuration is a $c$-element set
$M\subseteq\cT$ with $T_j\subseteq M$ while job $j$ is processed.  For a
job sequence $\pi=(\pi_1,\ldots,\pi_n)$ and a covering configuration sequence
$\cM=(M_1,\ldots,M_n)$, the number of switches is
\[
  z(\cM)=\sum_{i=2}^n |M_i\setminus M_{i-1}|.
\]
The initial loading is not charged.  Minimizing over covering histories gives
$z(\pi)$, and minimizing $z(\pi)$ over all job orders gives the global optimum.
This is the classical model of Tang and Denardo~\cite{TangDenardo1988}; their
keep-tool-needed-soonest (KTNS) policy optimizes the loading decisions when the
job sequence is fixed.

We use exactly the insertion search and the rule of
Locatelli, C\^ot\'e, and Coelho~\cite{LocatelliCoteCoelho2026}.  Its root is an
ordered orientation of a pair maximizing
\[
  |T_a\setminus T_b|+|T_b\setminus T_a|.
\]
At a partial sequence $\sigma=(\sigma_1,\ldots,\sigma_\ell)$, the next
unsequenced job maximizes
\[
  \min_{1\leq i\leq\ell}|T_{\sigma_i}\cup T_j|,
\]
and inserting that job in every position produces the children.  Reversal
symmetry retains only one orientation.  The almost-exact rule fathoms
$\sigma$ if, in a KTNS-optimal history $(M_1,\ldots,M_\ell)$, some adjacent
pair satisfies
\[
  T_{\sigma_i}\subseteq M_{i+1}
  \quad\text{and}\quad
  \sigma_i>\sigma_{i+1}.                                      \tag{1}
\]
A proper node is \emph{unsafe} when it is fathomed by (1) but some full order
obtained by inserting all remaining jobs into it has globally optimal cost.
The cited paper gives an unsafe event with five jobs.  The following theorem
shows that five is not a lower bound; it does not claim that every optimal
branch is lost or that the heuristic returns a suboptimal solution.

\begin{theorem}\label{thm:counterexample}
The assertion that no unsafe node exists in any instance with at most four
jobs, at most four tools, and
$\max_j|T_j|\leq c<m$ is false.
\end{theorem}

\begin{proof}
Consider
\[
 n=m=4,\qquad c=3,
 \qquad
 T_1=\{1\},\quad T_2=\{1,2\},\quad
 T_3=\{1,3,4\},\quad T_4=\{2,3,4\}.                         \tag{2}
\]
All four tools occur, every requirement is nonempty, and
$\max_j|T_j|=3=c<4=m$, so (2) is admissible.

For the unordered job pairs
$(1,2),(1,3),(1,4),(2,3),(2,4),(3,4)$, the symmetric-difference scores are,
respectively,
\[
  1,\ 2,\ 4,\ 3,\ 3,\ 2.                                   \tag{3}
\]
Thus $\{1,4\}$ is the unique root pair.  Relative to either of its
orientations, the remaining jobs have selection scores
\[
 \min_{i\in\{1,4\}}|T_i\cup T_2|=2,
 \qquad
 \min_{i\in\{1,4\}}|T_i\cup T_3|=3.                         \tag{4}
\]
Hence job $3$ is selected uniquely.  If the root orientation is $(1,4)$,
front insertion generates the proper node
$\sigma_+=(3,1,4)$.  If the orientation is $(4,1)$, middle insertion generates
$\sigma_-=(4,3,1)$.  Consequently the unspecified choice of root orientation
cannot remove the construction.

Write
\[
  A=\{1,3,4\},\qquad B=\{2,3,4\}.
\]
The unique KTNS-optimal magazine history for $\sigma_+$ is $(A,A,B)$.
Indeed, jobs $3$ and $4$ force $A$ and $B$, respectively, and retaining $A$
at the intervening job $1$ gives the unique one-switch history.  At the first
adjacent pair, $T_3=A\subseteq A=M_2$ and $3>1$, so (1) fathoms
$\sigma_+$.  Similarly, the unique KTNS-optimal history for $\sigma_-$ is
$(B,A,A)$.  At its second adjacent pair,
$T_3=A\subseteq A=M_3$ and $3>1$, so (1) fathoms $\sigma_-$.  Thus neither
event depends on a KTNS filler or eviction tie.

The only unsequenced job is job $2$.  Its requirement $\{1,2\}$ is contained
in neither $A$ nor $B$.  Therefore the completion-pruning test in Section~6.2
of~\cite{LocatelliCoteCoelho2026}, which requires every unsequenced job to be
covered by a node configuration, does not apply at either node.

It remains to prove global optimality of a completion.  Jobs $3$ and $4$
force the distinct full configurations $A$ and $B$.  Passing between them
already costs one switch.  If a complete history had total cost at most one,
all its other configurations would have to equal $A$ or $B$; otherwise an
additional transition would be charged.  Neither $A$ nor $B$ covers job $2$,
a contradiction.  Hence the global optimum is at least two.

Deleting job $2$ from $(2,3,1,4)$ leaves $\sigma_+$.  The covering history
\[
  (\{1,2,3\},A,A,B)
\]
has cost two.  Likewise, deleting job $2$ from $(2,4,3,1)$ leaves
$\sigma_-$, and
\[
  (\{1,2,3\},B,A,A)
\]
also has cost two.  Both full orders are therefore globally optimal
completions of their respective fathomed nodes.  Under either root
orientation, an unsafe node exists, proving the theorem.
\end{proof}

\section{Scope and reproducibility}

The five-job construction in Section~6.3 of
Locatelli, C\^ot\'e, and Coelho~\cite{LocatelliCoteCoelho2026} establishes the
same type of unsafe event, while the contribution here is the smaller explicit
instance (2), together with orientation robustness, unique selection scores,
forced KTNS histories, and the nonapplication of completion pruning.  We make
no assertion here about whether an unsafe instance exists with three jobs, or
about the final incumbent returned by the full heuristic.

The accompanying verifier reconstructs all size-three magazine histories,
evaluates all $24$ complete job orders in exact integer arithmetic, and checks
(2)--(4), both orientation cases, both fathoming triggers, and the failure of
completion pruning.  It returns global optimum two.  This computation is not
needed for the proof above.  For methodological comparison, Zeng~\cite{Zeng2026Whiskered}
likewise uses exhaustive exact arithmetic together with a short independent
verifier in a finite combinatorial setting; the present theorem does not depend
on that paper.

\paragraph{Computational and AI assistance.}
OpenAI Codex assisted in finding and checking the counterexample, preparing an
exact Python replay, checking cited source passages, and drafting the
manuscript.  The replay uses only finite enumeration and exact integer set
operations.  The mathematical claim rests on the explicit certificate and
proof above, not on software output or probabilistic computation.

\bibliographystyle{plain}
\bibliography{references}

\end{document}
